\chapter{DishBrain}
\chaptermark{DishBrain}
\label{sec:dishbrain}

\section{培养皿里的网络打网球}

DishBrain 是作者给一个试验台起的名字：在这个试验台上，培养在带电极芯片上的活神经元在玩单球拍的简化网球~\cite{Kagan2022}。这是用活神经元搭成的机器中讨论最多的一项实验（这类机器的综述见第~\ref{sec:livingmachine}~章），对本书来说也格外合用。这里的小网络整个都摸得着：每个输入由实验者设定，每个输出都被记录下来。前面各章的所有问题，即信号如何编码、谁来教、探针能看到什么，都可以拿来问一个既没有眼睛、没有肌肉、也没有 \nt{DOPA} 神经元的网络。我们像工程师剖析别人的试验台那样来剖析这项实验：电路、输入、输出、老师、结果、争议。

\section{试验台}

神经元取自小鼠胚胎皮层，每块芯片约 $800\,000$ 个细胞；或由人类干细胞培养而来，每块芯片约一百万个。它们在芯片上生长几周，彼此长入对方，并开始自发地发放脉冲和簇放电。在它们下方 $8$~mm$^2$ 的区域上有 $26\,000$ 个铂电极；最多可同时记录其中的 $1024$ 个，而实际用于刺激的只有八个。

培养物周围闭合着一个回路（图~\ref{fig:dishloop}）。计算机模拟游戏：球飞行、在墙上反弹，球拍上下移动。球的位置被转换成八个“感觉”电极上的电流。两个“运动”区的脉冲按 $10\ms$ 的窗口计数，用来移动球拍。每次击中或未击中后，培养物还会再收到一个刺激，即反馈。$10\ms$ 窗口是试验台计算机的时钟节拍，而不是培养物的：网络本身仍像本书中所有网络一样异步工作。

\begin{figure}[t]
\centering
\begin{tikzpicture}[>=Stealth,
  blk/.style={draw,very thick,rounded corners=3pt,align=center,font=\small,minimum height=1.2cm},
  lab/.style={font=\scriptsize,align=center}]
\node[blk,fill=blue!6,text width=3.4cm] (C) at (0,0) {芯片上的培养物\\\scriptsize $\sim10^6$ 个神经元，$26\,000$ 个电极};
\node[blk,fill=orange!10,text width=3.0cm] (G) at (7.2,0) {计算机：\\网球游戏};
\draw[->,very thick,blue!60!black] (G.north west) to[bend right=25] node[lab,above] {球：\emph{位置}——8 个电极中的哪一个，\\\emph{频率}——4--40 Hz} (C.north east);
\draw[->,very thick,red!70!black] (C.south east) to[bend right=25] node[lab,below] {球拍：“上”区和“下”区\\在 $10\ms$ 内的脉冲计数} (G.south west);
\draw[->,thick,densely dashed,green!45!black] (G.east) -- ++(0.9,0) |- ++(0,-2.6) -| node[lab,pos=0.25,below] {接住：可预测刺激；未接住：不可预测刺激} (C.south);
\end{tikzpicture}
\caption{DishBrain 回路。蓝色箭头是输入：球的位置用位置编码和频率编码表示（第~\ref{sec:code}~章）。红色箭头是输出：两个区的计数之差移动球拍。绿色虚线是每个回合之后的反馈，也是试验台上唯一的“老师”。}
\label{fig:dishloop}
\end{figure}

\section{输入与输出}

\emph{输入}同时用了第~\ref{sec:code}~章的两种编码。刺激八个电极中的哪一个，表示球在什么高度：这是位置编码，和第~\ref{sec:sensors}~章的传感器一样。刺激脉冲有多密，表示球有多远：球在远端墙边时每秒 $4$ 个，离球拍所在的墙越近线性增多，直到每秒 $40$ 个。这是频率编码；球的刺激脉冲幅度为 $75$~mV。培养物事先并不“认识”这两种编码中的任何一种：编码由实验者规定，网络必须学会读懂它。

\emph{输出}的读取方式与第~\ref{sec:debugging}~章的接口相同。一个运动区的脉冲把球拍往上推，另一个往下推；最简单的写法是，每个窗口内
\begin{empheq}[box=\keybox]{equation}
\Delta p \;=\; c\,\bigl(n_\uparrow - n_\downarrow\bigr),
\label{eq:dishreadout}
\end{empheq}
其中 $n_\uparrow$、$n_\downarrow$ 是两个区在 $10\ms$ 内的脉冲数，$c$ 是试验台的系数。这是第~\ref{sec:glutcomputer}~章的双线编码：运动方向不是由信号的符号携带，而是由两根导线中哪一根更响来携带。

来看看这种输出的噪声。如果一个区每秒总共发出 $R$ 个脉冲，那么在 $T=10\ms$ 的窗口里平均有 $RT$ 个，按~\eqref{eq:rateerror}，相对离散度为 $1/\sqrt{RT}$。$R=1000$ 次/秒时每个窗口约为百分之 $30$；$R=100$ 时超过百分之百。球拍不可避免地会抖动，网络只能用一种方式学习：把计数之差的\emph{平均值}往正确方向推。这正是限制第~\ref{sec:debugging}~章中任何接口的同一套规律，只不过现在出现在回路的两侧。

\section{没有多巴胺的老师}

试验台最有意思的部分是老师。皮层神经元培养物里没有 \nt{DOPA} 神经元，试验台也不发送任何“比预期好”的信号。反馈是另一种形式。网络接住了球，所有八个感觉电极同时收到一个短暂的\emph{可预测}刺激：$75$~mV 的脉冲，每秒 $100$ 个，持续 $0.1\,\text{s}$。没接住，就是四秒钟作者称为\emph{不可预测}的刺激：$150$~mV 的脉冲，每秒 $5$ 个，在随机时刻施加到随机选择的电极上；随后是四秒钟不刺激的停顿，球重新发出~\cite{Kagan2022}。

作者的出发点是一个假说：网络的行为是为了让自己的输入变得可预测~\cite{Friston2010}。对工程师来说含义很简单：培养物摆脱四秒随机噪声的办法只有一个，就是接住球。同样的思想我们在第~\ref{sec:code}~章的节省脉冲和第~\ref{sec:recognition}~章的画面预测中都见过。

但这里需要的真是可预测性吗？还可能有更粗糙的机制。假设未接住后的不可预测刺激只是把网络中近期的变化\emph{打乱}，而接住后的平静刺激不去碰它们。用一个向量 $\boldsymbol\psi$ 描述网络的设置，设在设置 $\boldsymbol\psi$ 下网络未接住的概率为 $P_{\text{miss}}(\boldsymbol\psi)$。如果扰动是对称的，即从 $\boldsymbol\psi$ 到 $\boldsymbol\psi'$ 与反方向的可能性相同，那么在稳态下，离开每个设置的流量与进入的流量相平衡，网络停留在设置 $\boldsymbol\psi$ 的时间比例为
\begin{empheq}[box=\keybox]{equation}
\rho(\boldsymbol\psi) \;\propto\; \frac{1}{P_{\text{miss}}(\boldsymbol\psi)} .
\label{eq:dishdensity}
\end{empheq}
未接住的次数少十倍的设置，保持的时间就长十倍。没有人告诉网络该往哪个方向改权重：好的设置只是更少被破坏。

我们在一个模型上检验了这一点，把培养物简化为两个数：球到达高度 $y$ 时，球拍到达 $a\,y+b$，偏差小于 $0.1$ 就算接住。未接住后两个数都受到一次随机扰动，接住后保持不变。在 $2000$ 个回合中，接住的比例从前两百个回合的 $0.21$ 上升到最后五百个回合的 $0.38$（四十个“培养物”，图~\ref{fig:dishmodel}）。如果在\emph{每个}回合后都给扰动，反馈中就没有信息，比例停留在 $0.12$ 左右；如果完全没有扰动，比例保持在 $0.19$。试验台上的实验与此一致：一次会话内的显著提升只出现在完整反馈下。以安静代替刺激时，培养物到会话结束时仍比没有反馈时打得好，但效果较小~\cite{Kagan2022}；需要指出，在这种条件下，未接住后也跟着较长的停顿，接住后则较短，因此两种结局之间的差别并未完全消失。

\begin{figure}[t]
\centering
\begin{tikzpicture}[x=1cm,y=5cm]
\draw[->,gray!70!black] (0,0) -- (10.6,0);
\draw[->,gray!70!black] (0,0) -- (0,0.5) node[above,font=\small,black] {接住比例};
\foreach \v in {0.1,0.2,0.3,0.4}{\draw[gray!60!black] (-0.06,\v) -- (0.06,\v); \node[left,font=\scriptsize] at (0,\v) {\v};}
\foreach \x/\f/\l/\lab in {1.8/0.21/0.38/{未接住后\\给噪声},5.2/0.13/0.11/{每个回合后\\都给噪声},8.6/0.19/0.19/{无噪声}}{
  \fill[blue!35] (\x-0.8,0) rectangle (\x-0.05,\f);
  \fill[red!60!black] (\x+0.05,0) rectangle (\x+0.8,\l);
  \node[below,font=\scriptsize,align=center] at (\x,-0.01) {\lab};
  \node[above,font=\scriptsize] at (\x-0.425,\f) {\f};
  \node[above,font=\scriptsize] at (\x+0.425,\l) {\l};}
\fill[blue!35] (7.4,0.47) rectangle (7.7,0.495); \node[right,font=\scriptsize] at (7.75,0.482) {前 200 个};
\fill[red!60!black] (7.4,0.41) rectangle (7.7,0.435); \node[right,font=\scriptsize] at (7.75,0.422) {后 500 个};
\end{tikzpicture}
\caption{“未接住即打乱”模型（\texttt{dishbrain\_model.py}）：$2000$ 个回合开始和结束时的接住比例，四十个模型培养物的平均值。只有当扰动跟在未接住之后、而不跟在接住之后时，学习才会发生，这与实验一致：一次会话内的显著提升只出现在完整反馈下。}
\label{fig:dishmodel}
\end{figure}

\begin{keyblock}
\begin{proposition}[打乱式老师]
\label{prop:dishscramble}
如果失败会晃动网络、成功则让它保持原样，网络就会自行漂向更少出错的设置：停留在某个设置的时间与失败概率成反比~\eqref{eq:dishdensity}。这种学习既不需要奖励信号，也不需要其符号，也不需要可预测性本身，只需要成功后和失败后的反馈有所不同。培养物的行为变化已测量；培养皿里起作用的是哪种机制，是输入预测还是打乱，尚未确定。
\end{proposition}
\end{keyblock}

与第~\ref{sec:dofa}~章比较一下。那里噪声也用于搜索，但多巴胺带有符号：误差 $\delta$ 指明权重该往哪个方向移动，步子沿梯度走~\eqref{eq:perturbation}。这里没有符号，只有“保留”或“晃动”。这种搜索更粗糙、更慢，但对网络的要求更少：不需要 \nt{DOPA} 神经元，不需要评论家，也不需要延续数秒的资格迹。

\section{测到了什么，争的是什么}

每局游戏持续 $20$ 分钟；比较的是前 $5$ 分钟和后 $15$ 分钟。在有完整反馈的培养物中，连续击球的回合变长，发球即丢的情况变少；在没有细胞、没有游戏以及由计算机控制“球拍”的对照中，都没有类似现象。人类细胞培养物平均能比小鼠培养物坚持更久。这一改进在统计上是可靠的，涉及九个小鼠培养物和十一个人类培养物，每组一百多局游戏，但幅度不大：网络并没有成为好球手，只是比随机稍好一点。它在一次会话内可靠；在几天里跨会话的稳定学习，作者并没有得到~\cite{Kagan2022}。同一小组的后续工作发现，游戏时培养物的活动会趋近\emph{临界}状态，即衰减与雪崩之间的边界，也就是第~\ref{sec:gaba}~章中那种处在边缘的状态，而且越接近临界，游戏打得越好~\cite{Habibollahi2023}。

引起主要争议的不是数字，而是措辞。文章标题宣称培养皿中的神经元“表现出感知能力”（sentience），一组研究者对此提出反驳：闭环中的行为变化还不是智能，也不是感受，结论应当表述得更谨慎~\cite{Balci2023}。从工程角度看，有两个问题悬而未决，都关乎机制。第一，网络是在学习，即它的权重长期改变了，还是反馈只是移动了它的当前状态？第二，需要的真是可预测性，还是只要对成功和失败的响应有任何差别就够了，如命题~\ref{prop:dishscramble} 所述？在每个电极都可访问的试验台上，这两个问题都可以回答，这或许正是它最大的价值。

\section{试验台规格：如何复现}
\label{sec:dishspec}

下面汇总了重新搭建这样一个试验台所需的全部内容：设备、回路、输入编码、输出读取、反馈和实验方案。每个参数都注明了来源。“论文”表示数值取自文献~\cite{Kagan2022} 的正文。“自选”表示论文中没有给出，复现时需要自行设定；我们给出默认值，并说明如何检验。

\subsection*{设备与培养物}

\begin{center}
\footnotesize
\setlength{\tabcolsep}{4pt}
\renewcommand{\arraystretch}{1.15}
\begin{tabular}{>{\raggedright}p{4.0cm}>{\raggedright}p{6.9cm}>{\raggedright\arraybackslash}p{1.6cm}}
\toprule
参数 & 数值 & 来源 \\
\midrule
芯片 & MaxOne 阵列：$8$~mm$^2$ 区域上 $26\,000$ 个铂电极 & 论文 \\
记录 & 最多同时 $1024$ 个通道，每秒 $20\,000$ 个采样 & 论文 \\
刺激 & 技术上最多 $32$ 个电极，实验中用 $8$ 个独立电极 & 论文 \\
细胞 & 小鼠胚胎皮层；由干细胞获得的人类神经元（两种培养方法）；HEK293T 细胞（非神经元）作对照 & 论文 \\
接种 & 每块芯片约 $10^6$ 个细胞 & 论文 \\
实验时的培养龄 & 小鼠：从 $\sim10$ 天起；人类：视培养方法为 $14$--$24$ 天或 $\sim80$ 天 & 论文 \\
\bottomrule
\end{tabular}
\end{center}

\subsection*{回路与时间}

回路以 $10\ms$（$200$ 个采样）为窗口运行：每个窗口内统计运动区电极上的脉冲，更新游戏，并发出下一次刺激。从培养物中的脉冲到响应刺激的延迟约 $5\ms$。游戏期间，每个电极的信号经过截止频率 $100$~Hz 的二阶贝塞尔高通滤波器；信号的绝对值再经 $1$~Hz 的一阶贝塞尔低通滤波器平滑，脉冲阈值与这个平滑后的绝对值成正比。论文中六倍背景标准差的阈值用于单独的活动测量，而不是游戏。为了不把刺激误算为培养物的响应，当同时出现超过 $15$ 个大于 $75$~mV 的大脉冲时，所有信号都被屏蔽。以上均来自论文。

\subsection*{球的编码}

\begin{center}
\footnotesize
\setlength{\tabcolsep}{4pt}
\renewcommand{\arraystretch}{1.15}
\begin{tabular}{>{\raggedright}p{3.4cm}>{\raggedright}p{7.5cm}>{\raggedright\arraybackslash}p{1.6cm}}
\toprule
参数 & 数值 & 来源 \\
\midrule
感觉区 & $8$ 个电极，排列使相邻电极表示球的相邻位置 & 论文 \\
位置编码 & 电极编号 $k=1,\dots,8$ 表示球相对于球拍中心的位置 & 论文 \\
哪个坐标 & 球的竖直位置；复现时取相对于球拍的位置，$y-p\in[-\tfrac12,\tfrac12]$ 时 $k=\lceil 8(y-p+\tfrac12)\rceil$ & 论文；公式为自选 \\
频率编码 & $4$ 到 $40$~Hz 的刺激频率表示到球拍的距离 & 论文 \\
标度方向 & 球在对面墙边时 $4$~Hz，线性增加到球在球拍一侧墙边时的 $40$~Hz：$f=4+36\,(1-x/L)$，$x$ 为到球拍一侧墙的距离，$L$ 为场地宽度 & 论文 \\
脉冲波形 & 短的矩形双相脉冲：先正相，后负相 & 论文 \\
幅度 & $75$~mV；选定这一幅度是为了不强迫受抑制的细胞发放 & 论文 \\
相位时长 & 未给出；采用刺激系统的标准设置，实验期间不改变 & 自选 \\
\bottomrule
\end{tabular}
\end{center}

\subsection*{读取球拍}

论文中有两个运动区：第一个区的脉冲让球拍上移，第二个区的让它下移；两个区的贡献经过加权，以补偿细胞密度和活动水平的差异~\cite{Kagan2022}。确切公式没有给出。复现时我们采用规则~\eqref{eq:dishreadout}，每个 $10\ms$ 窗口 $\Delta p=c\,(n_\uparrow-n_\downarrow)$，用权重按静息平均活动把两个区拉平（自选）：$n_\uparrow$ 和 $n_\downarrow$ 各自除以游戏前测得的本区每窗口平均计数。系数 $c$ 取为使一个平均计数的差每窗口把球拍移动场地高度的 $1\%$（自选）；这样，一个区持续占优时，球拍 $1$~s 就能走完整个场地。

论文正文没有用坐标给出各区在芯片上的位置，只有一张示意图。复现时只需三组互不重叠的电极：中间八个感觉电极，两侧各一组记录电极，一组为“上”，一组为“下”（自选）。

\subsection*{游戏}

论文没有给出场地尺寸、球拍长度和球速；只说球以恒定速度飞行并在墙上反弹。复现时（自选）：场地 $1\times1$，球拍在右侧墙上，长 $0.2$（$|y-p|<0.1$ 即击中，与模型~\texttt{models/dishbrain\_model.py} 一致），球 $1$~s 飞越场地。一个回合中一次都没接住就未接住，记为“ace”；连续击球超过三次的回合记为“长回合”（论文）。

\subsection*{反馈}

\begin{center}
\footnotesize
\setlength{\tabcolsep}{4pt}
\renewcommand{\arraystretch}{1.15}
\begin{tabular}{>{\raggedright}p{2.6cm}>{\raggedright}p{8.3cm}>{\raggedright\arraybackslash}p{1.6cm}}
\toprule
事件 & 施加的刺激 & 来源 \\
\midrule
击中 & 可预测刺激：全部 $8$ 个感觉电极同时，$75$~mV，$100$~Hz，$100\ms$；其间暂停通常的球位置刺激 & 论文 \\
未击中 & 不可预测刺激：$150$~mV，$5$~Hz，$4$~s，在随机时刻施加到从 $8$ 个电极中随机选出的电极上；随后 $4$~s 无刺激，球沿随机方向重新发出 & 论文 \\
随机规律 & 未给出；复现时把脉冲时刻取为平均频率 $5$~Hz 的泊松流 & 自选 \\
\bottomrule
\end{tabular}
\end{center}

\subsection*{实验方案与对照}

一次会话持续 $20$~min；比较前 $5$~min 和后 $15$~min（论文）。三项结果指标：平均回合长度、ace 比例和长回合比例。实验条件（论文）：
\begin{itemize}
\item \emph{刺激}：完整回路，如上所述；
\item \emph{安静}：以同样长的无任何刺激时段代替击中或未击中刺激，然后球沿随机方向发出；
\item \emph{无反馈}：未接住后球直接反弹继续飞行，球位置刺激继续进行；
\item \emph{静息}：游戏进行、球拍随活动移动，但完全没有刺激；
\item \emph{对照}：只有培养基的芯片；HEK293T 细胞；“计算机里的培养物”，即用仿真代替细胞。
\end{itemize}
主系列的样本量：小鼠培养物 $9$ 个（$101$ 次会话），人类培养物 $11$ 个（$138$），只有培养基的芯片 $6$ 块（$80$），静息培养物 $20$ 个（$42$），仿真 $3$ 个（$38$），共 $399$ 次会话。反馈系列：$3$ 天、每天 $3$ 次会话，共 $486$ 次会话，条件为“刺激”“安静”和“无反馈”（$n=27$、$17$ 和 $15$）。平均回合长度从前 $5$ 分钟到后 $15$ 分钟的增长只出现在活的培养物中。在反馈系列中，随时间的显著增长只出现在“刺激”条件下；会话结束时，“安静”条件仍优于“无反馈”，但效应较小；而在三天里跨会话的稳定学习并未出现~\cite{Kagan2022}。

\subsection*{试验台循环的步骤}

每隔 $10\ms$，试验台计算机执行以下操作。
\begin{enumerate}
\item 读取上一窗口内两个运动区的脉冲数 $n_\uparrow$、$n_\downarrow$，并用各区静息平均计数归一化。
\item 把球拍移动 $\Delta p=c\,(n_\uparrow-n_\downarrow)$，并限制在场地边界内。
\item 把球移动一步；在上、下、左三面墙上反弹。
\item 如果球到达右墙：若 $|y-p|<0.1$ 则为击中，回合计数加一，施加击中刺激（$100\ms$）；否则为未击中，记录该回合，施加未击中刺激（$4$~s），再停顿 $4$~s，球重新发出。
\item 否则，按球的竖直位置选择电极 $k$，按到球拍一侧墙的距离选择频率 $f$，在下一个窗口之前向电极 $k$ 施加频率为 $f$、幅度 $75$~mV 的双相脉冲。
\end{enumerate}
这些足以搭建一个与已发表试验台兼容的装置，并检验本章的核心问题：教会培养物的是可预测性，还是仅仅是对击中与未击中响应之间的差别。为此，应在论文的三个条件之外加上论文中没有的第四个：未接住后施加与不可预测刺激时长和能量相同的\emph{可预测}刺激。如果培养物这样也能学会，关键就在于差别，而不在于可预测性。


\section{小结}

\begin{table}[t]
\centering
\footnotesize
\setlength{\tabcolsep}{4pt}
\renewcommand{\arraystretch}{1.15}
\begin{tabular}{>{\raggedright}p{2.1cm}>{\raggedright}p{4.2cm}>{\raggedright}p{1.9cm}>{\raggedright\arraybackslash}p{4.6cm}}
\toprule
试验台部件 & 结构 & 本书章节 & 已知情况 \\
\midrule
输入 & 位置（8 个电极中的哪一个）和频率（4--40 Hz） & \ref{sec:code}, \ref{sec:sensors} & 由实验者规定 \\
输出 & 两个区在 $10\ms$ 内的计数之差~\eqref{eq:dishreadout} & \ref{sec:glutcomputer}, \ref{sec:debugging} & 由实验者规定；噪声按~\eqref{eq:rateerror} \\
老师 & 接住后给可预测刺激，未接住后给不可预测刺激；没有多巴胺 & \ref{sec:dofa} & 一次会话内的显著提升只出现在完整反馈下；安静条件下效应较小；跨会话不稳定 \\
机制 & 输入预测或未接住即打乱~\eqref{eq:dishdensity} & \ref{sec:dofa}, \ref{sec:recognition} & 未确定；两者都能解释结果 \\
网络状态 & 游戏时趋近临界状态 & \ref{sec:gaba} & 与游戏水平的关系已测量 \\
“感知能力” & 作者的论断 & --- & 未被证明；有争议 \\
\bottomrule
\end{tabular}
\caption{工程师眼中的 DishBrain。}
\label{tab:dishbrain}
\end{table}

DishBrain 是最简形式的活神经元机器：输入用编码给定，输出靠计数读取，老师就是对成功与失败的响应之差（表~\ref{tab:dishbrain}）。一个没有眼睛、肌肉和多巴胺的网络稍微学会了打球，单凭这一点，它就成了检验本书思想的试验台。无论最后的机制是预测、打乱还是别的什么，答案都会按第~\ref{sec:debugging}~章建议的方式找到：用探针、测试信号和关闭部件，在一台终于能整个看清的机器上。

\section{习题}

\begin{exercise}[\lvlA]
\label{ex:22:1}
\emph{计数要偏移多少。}两个区合计每秒发出 $R_\uparrow+R_\downarrow=2000$ 个脉冲，均为泊松流。(a)~$R=1000$ 和 $R=100$ 时，一个区在 $10\ms$ 内计数的相对离散度是多少？(b)~球飞行 $1$~s 期间位移~\eqref{eq:dishreadout} 累加。$R_\uparrow-R_\downarrow$ 至少多大，平均位移才是离散度的两倍？
\end{exercise}

\begin{exercise}[\lvlA]
\label{ex:22:2}
\emph{要多少回合才能看出学习。}各回合独立，接住比例服从二项分布。两个时段中各需多少回合，才能使接住比例的增长超过差值的两个标准差：(a)~从 $0.20$ 到 $0.25$；(b)~如模型中那样从 $0.21$ 到 $0.38$？
\end{exercise}

\begin{exercise}[\lvlB]
\label{ex:22:3}
\emph{推导~\eqref{eq:dishdensity}。}设置 $\boldsymbol\psi$ 取值于有限集合；未接住时（概率 $P(\boldsymbol\psi)>0$），网络以对称概率 $K(\boldsymbol\psi,\boldsymbol\psi')$ 转到 $\boldsymbol\psi'$，接住时保持不变。(a)~证明 $\rho\propto1/P$ 是平稳分布。(b)~两个设置分别为 $P=0.9$ 和 $P=0.1$ 时，未接住的比例是多少？
\end{exercise}

\begin{exercise}[\lvlB]
\label{ex:22:4}
\emph{模型的吸收区。}球拍到达 $p=\min\bigl(1,\max(0,ay+b)\bigr)$，球到达高度 $y\sim U(0,1)$，$|p-y|<h=0.1$ 即接住。(a)~证明当 $|b|<h$ 且 $|a-1+b|<h$ 时网络接住所有球，并求该区域的面积。(b)~数值求 $a\sim U(-1,1)$、$b\sim U(0,1)$ 时的平均接住比例。
\end{exercise}

\begin{exercise}[\lvlB]
\label{ex:22:5}
\emph{建模：给设置围上栅栏。}在 \texttt{dishbrain\_model.py} 的副本中，让 $200$ 个培养物各打 $20\,000$ 个回合。最后 $500$ 个回合的接住比例是多少？如果每次扰动后把设置裁剪到矩形 $a\in[-1,2]$、$b\in[-1,1]$ 内呢？用习题~\ref{ex:22:3} 解释差别。
\end{exercise}

\begin{exercise}[\lvlB]
\label{ex:22:6}
\emph{设计输入编码。}高度误差小于 $h=0.1$ 即接住；网络在 $T=0.25$~s 内读取输入，频率不超过 $r$ 时约可区分 $\sqrt{rT}$ 个等级，刺激不超过每秒 $40$ 个脉冲。(a)~八个电极上的位置编码够用吗？(b)~能否用单个电极上的频率来传递高度？
\end{exercise}

\begin{exercise}[\lvlC]
\label{ex:22:7}
\emph{研究：预测还是打乱？}把模型推广到 $d$ 个可调数（$p=\sum_{i<d}a_iy^i$）。在打乱规则下和在带误差符号的规则~\eqref{eq:perturbation} 下，达到接住比例 $0.5$ 所需的回合数如何随 $d$ 增长？试验台上什么实验能区分这两个假说？
\end{exercise}
